\documentclass{article}

\addtolength{\topmargin}{-1in}
\addtolength{\textheight}{1.75in}
\title{Research Statement}
\author{Gautam Kamath}
\date{}
\begin{document}
\maketitle
\pagenumbering{gobble}
My research interests lie in the field of theoretical computer science,
specifically in the areas of algorithmic game theory, stochastic processes,
and applied probability.

Residential segregation is a phenomenon which can be observed in many
major metropolitan areas in the United States, including New York City and Chicago.
Social scientists are often interested in studying and understanding the dynamics
through which segregation occurs.
In 1969, economist Thomas Schelling introduced a mathematical model to mimic this
process.
In Schelling's segregation model, individuals of two types live on a ring.
Individuals are unhappy if their type is a minority within their local neighborhood.
At each time step, two unhappy individuals swap positions.
This process continues until all individuals are happy.

Schelling ran simulations by hand and reported the results.
While individuals in the model would be happy living in a fairly integrated state,
the resulting configurations exhibited significant segregation:
individuals ended up in long monochromatic runs, made up entirely of individuals of the same type.

While many were interested by this model, most research results were experimental, rather than analytic.
The first major analysis was by Young in 2001, where he studied
the effect of perturbations on the model.
His showed that, in the limit as perturbations tended to $0$, complete segregation was inevitable.

In joint work with Brandt, Kleinberg and Immorlica in 2012, we studied the
unperturbed model, and found a very different outcome.
Our results showed two things:
First, that segregation is local, rather than global.
By this, we mean that the degree of segregation depends solely on the size of an individual's neighborhood,
and is independent of the size of the world.
We also showed that segregation is modest.
The degree of segregation depends linearly on the size of an individual's neighborhood, rather than exponentially.

Our technique involves splitting the ring into small segments 
and studying the probability that a stable configuration forms in a segment.
We show that monochromatic runs of opposite types occur fairly often in the ring,
thus limiting the length of any single run.
In order to study a small portion of the ring in isolation,
we use Wormald's technique of differential equations to provide guarantees
about the global state of the ring.

My present work focuses on a classic learning problem:
the Gaussian mixture model.
We are given black-box access to a distribution $F$
with density function $F = \sum_{i=1}^k w_iF_i$, where $\sum_{i=1}^k w_i = 1$
and $F_i \sim \mathcal{N}(\mu_i,\sigma^2_i)$.
The goal is to learn a distribution $F'$ with low distance to $F$.

In a series of papers by Kalai, Moitra, and Valiant, it is shown
that a mixture of Gaussian can be learned with a time and data
requirement polynomial in the dimension and the inverse of the desired accuracy,
and exponential in the number of Gaussians.
Their technique involves the use of random projections down to one dimension and applying
the method of moments to each subproblem.
However, they emphasize that they favor clarity of proof and exposition over
optimization of runtime.
As such, the exponents in the runtime are much greater than one could tolerate
in practice.
In joint work with Daskalakis, we seek to find an alternate method of
learning a mixture of Gaussians, with a time requirement that could
result in practical use.
\end{document}
